% Chapter — Reference frames (FR-3). Follows the mandatory FR-29
% six-section template: purpose; assumptions and domain of validity;
% derivation; discretization and implementation; validation evidence;
% references. The equation labels eq:frames:fw, eq:frames:nut00b,
% eq:frames:xy, eq:frames:s06, eq:frames:c2i, eq:frames:era,
% eq:frames:chain, eq:frames:moonpa, and eq:frames:mars are echoed
% verbatim in comments in cpp/src/frames.cpp and
% cpp/include/star/frames.hpp (FR-29 equation-label-to-code traceability);
% renaming them breaks that trace.
\chapter{Reference Frames}
\label{ch:frames}

\section{Purpose}
\label{sec:frames:purpose}

This chapter documents the reference-frame module (FR-3): the CIO-based
transformation between the GCRF ($I$ in
Table~\ref{tab:notation:frames}) and the rotating Earth-fixed frame ($E$)
using the IAU 2006/2000B precession-nutation model and the Earth rotation
angle; the lunar principal-axis frame ($L$) realized from the DE 3-1-3
libration Euler angles; and the Mars body-fixed frame ($M$) from the IAU
Working Group 2015 rotational elements. These constructions orient every
body-fixed model of later phases --- harmonic gravity fields, atmospheres,
ground-track and launch-site geometry --- and consume the time systems of
Chapter~\ref{ch:time} and the rotation kernel of
Chapter~\ref{ch:rotations}. The module is implemented in
\texttt{cpp/src/frames.cpp} with its interface in
\texttt{cpp/include/star/frames.hpp}.

\section{Assumptions and domain of validity}
\label{sec:frames:domain}

Assumptions:
\begin{enumerate}
  \item \textbf{IAU 2000B nutation suffices.} The 77-term IAU 2000B
    abridged series delivers a celestial pole within \qty{1}{\mas}
    ($\approx \qty{4.8e-9}{\radian}$) of the full IAU 2000A model over
    1995--2050~\cite{mccarthy2003}. A \qty{1}{\mas} pole error displaces an
    Earth-fixed position by $\sim \qty{3}{\centi\metre}$ at the surface
    --- far below the mission-analysis needs of this simulator. The free
    core nutation, unpredictable and of comparable sub-mas size, is
    unmodeled in both 2000A and 2000B.
  \item \textbf{Polar motion is neglected.} The produced ``Earth-fixed''
    frame is strictly the Terrestrial Intermediate Reference System: the
    final polar-motion rotation $\mat{W}(x_p, y_p, s')$ of the full IERS
    chain is omitted. Polar motion is observed, not modeled; its
    magnitude is bounded by $|x_p|, |y_p| \lesssim \qty{0.6}{\arcsecond}
    \approx \qty{2.9e-6}{\radian}$~\cite{iers2010}, i.e.\ up to
    $\sim \qty{18}{\metre}$ of surface displacement. EOP time-series
    ingestion is a PRD non-goal; the omission is exactly this bounded
    rotation.
  \item \textbf{Constant, user-suppliable $\Delta$UT1 (default 0).}
    UT1 is realized as $\mathrm{UT1} = \mathrm{UTC} + \Delta\mathrm{UT1}$
    with a constant offset. By construction of UTC,
    $|\mathrm{UT1} - \mathrm{UTC}| < \qty{0.9}{\second}$, so the default
    $\Delta\mathrm{UT1} = 0$ bounds the Earth-rotation error by
    \qty{0.9}{\second} of spin: $\qty{0.9}{\second} \times
    \qty{7.292e-5}{\radian\per\second} \approx \qty{6.6e-5}{\radian}$,
    $\sim \qty{420}{\metre}$ at the surface. A user-supplied constant
    appropriate to the mission epoch reduces the error to the UT1 drift
    accumulated across the mission window (of order milliseconds per day,
    i.e.\ decimetres per day at the surface).
  \item \textbf{Time plumbing.} The TT argument of the precession-nutation
    series comes from the two-part TAI epoch (Chapter~\ref{ch:time});
    the Mars elements use the D-6 truncated TDB. The TDB truncation error
    ($\leq \qty{30}{\micro\second}$, observed \qty{5.5e-6}{\second})
    enters Mars $W$ through the rotation rate as
    $\leq \qty{2.1e-9}{\radian}$ --- below the physical accuracy of the
    rotational-element model itself (mas level).
  \item \textbf{Moon PA assembly only.} The lunar construction is the
    exact 3-1-3 composition of supplied libration angles; the angles
    themselves come from the DE440 ephemeris evaluator (a separate
    module), and their physical accuracy is that of
    DE440~\cite{park2021}.
\end{enumerate}

Domain of validity: epochs from 1972 (the leap-table domain of
Chapter~\ref{ch:time}); model fidelity per assumption 1 holds over the
published 1995--2050 comparison span and degrades gracefully outside it
(the golden gates run to 2060 against the same model). The frame chain is
\emph{model-faithful} to IAU 2006/2000B at the $10^{-11}$ matrix-element
level everywhere in 2020--2060; its \emph{physical} fidelity to the true
Earth orientation is set by assumptions 2 and 3, which dominate every
other term.

Out-of-domain behavior, matching the code exactly: epochs before
1972-01-01 throw \texttt{std::domain\_error} from the time module (the
UT1 realization needs $\Delta AT$); epochs past the leap-table expiry use
the last tabulated offset with the Python-layer warning of
Chapter~\ref{ch:time}. No function of this module throws on its own.

\section{Derivation}
\label{sec:frames:derivation}

\subsection{The CIO-based chain}

The IERS Conventions (2010), Chapter~5~\cite{iers2010}, factor the
celestial-to-terrestrial transformation as $\mat{W}(t)\, \mat{R}_3(
\mathrm{ERA})\, \mat{Q}^{\mathsf T}(t)$, where $\mat{Q}(t)$ orients the
Celestial Intermediate Pole (CIP) and Origin (CIO) in the GCRS, ERA is the
Earth rotation angle, and $\mat{W}$ is polar motion. With $\mat{W}$
omitted (assumption 2) and this project's frame-transformation DCM
convention, the implemented chain is
\begin{equation}
  \dcm{I}{E} = \mat{R}_3\bigl(\mathrm{ERA}(\mathrm{UT1})\bigr)\,
               \dcm{I}{\mathrm{CIRS}},
  \label{eq:frames:chain}
\end{equation}
with $\dcm{I}{\mathrm{CIRS}}$ the celestial-to-intermediate matrix built
from the CIP coordinates $X, Y$ and the CIO locator $s$ below.

\subsection{IAU 2006 precession in Fukushima--Williams angles}

The IAU 2006 (P03) precession~\cite{capitaine2003} is parameterized by the
four Fukushima--Williams angles~\cite{fukushima2003, wallace2006}
$\bar\gamma$, $\bar\phi$, $\bar\psi$, $\epsilon_A$ --- polynomials in the
TT Julian-century argument $t$ of equation~\eqref{eq:time:ttcent}, with
frame bias absorbed into the polynomial constant terms. The
bias-precession-nutation (NPB) matrix folds the nutation components
directly into two of the angles~\cite{wallace2006}:
\begin{equation}
  \mat{M}_{\mathrm{NPB}} =
  \mat{R}_1\bigl(-(\epsilon_A + \Delta\epsilon)\bigr)\,
  \mat{R}_3\bigl(-(\bar\psi + \Delta\psi)\bigr)\,
  \mat{R}_1(\bar\phi)\,
  \mat{R}_3(\bar\gamma),
  \label{eq:frames:fw}
\end{equation}
using the elementary frame rotations of
equation~\eqref{eq:rotations:r1r2r3}. The polynomial coefficients (and the
IAU 2006 mean obliquity $\epsilon_A$) are tabulated in
\texttt{cpp/src/frames\_series.hpp}, transcribed from the ERFA 2.0.1
reference source~\cite{erfa201} implementing~\cite{capitaine2003,
wallace2006}.

\subsection{IAU 2000B nutation}

The abridged IAU 2000B model~\cite{mccarthy2003} evaluates 77 luni-solar
terms on the linear Delaunay arguments $l, l', F, D, \Omega$ (their
truncation to linear polynomials is part of the published model), plus
fixed offsets standing in for the truncated planetary terms:
\begin{equation}
  \Delta\psi = \sum_{k} \bigl[(A_k + A'_k t)\sin\theta_k
    + A''_k \cos\theta_k\bigr] + \Delta\psi_{\mathrm{pl}},
  \qquad
  \Delta\epsilon = \sum_{k} \bigl[(B_k + B'_k t)\cos\theta_k
    + B''_k \sin\theta_k\bigr] + \Delta\epsilon_{\mathrm{pl}},
  \label{eq:frames:nut00b}
\end{equation}
where $\theta_k$ is the integer combination of the five arguments for term
$k$ and the planetary offsets are $\Delta\psi_{\mathrm{pl}} =
\qty{-0.135}{\mas}$, $\Delta\epsilon_{\mathrm{pl}} = \qty{0.388}{\mas}$
(the Luzum values appropriate to applying frame bias, precession, and
nutation separately, as equation~\eqref{eq:frames:fw}
does~\cite{mccarthy2003, erfa201}).

\subsection{CIP coordinates, the CIO locator, and the intermediate matrix}

The CIP unit vector in the GCRS is the third row of the NPB matrix, so its
coordinates are read off directly:
\begin{equation}
  X = \bigl(\mat{M}_{\mathrm{NPB}}\bigr)_{31},
  \qquad
  Y = \bigl(\mat{M}_{\mathrm{NPB}}\bigr)_{32}.
  \label{eq:frames:xy}
\end{equation}
The CIO locator $s$ positions the Celestial Intermediate Origin on the
CIP equator; it is evaluated in the standard $s + XY/2$ form --- a
polynomial plus Poisson series in the eight fundamental luni-solar and
planetary arguments, per IERS Conventions (2010) Chapter~5 and the SOFA
S06 realization consistent with IAU 2006 precession~\cite{iers2010,
wallace2006}:
\begin{equation}
  s = -\frac{XY}{2}
      + \sum_{j=0}^{5} t^{j} \Bigl[ c_j + \sum_{k}
        \bigl( S_{jk} \sin\alpha_{jk} + C_{jk} \cos\alpha_{jk} \bigr)
        \Bigr],
  \label{eq:frames:s06}
\end{equation}
with the coefficients again tabulated in
\texttt{cpp/src/frames\_series.hpp}. Writing the CIP direction in
spherical form, $E = \operatorname{atan2}(Y, X)$ and
$d = \arctan\sqrt{(X^2 + Y^2)/(1 - X^2 - Y^2)}$, the
celestial-to-intermediate matrix is
\begin{equation}
  \dcm{I}{\mathrm{CIRS}} = \mat{R}_3\bigl(-(E + s)\bigr)\,
    \mat{R}_2(d)\, \mat{R}_3(E),
  \label{eq:frames:c2i}
\end{equation}
the standard construction of IERS Conventions (2010),
Chapter~5~\cite{iers2010}.

\subsection{Earth rotation angle}

The Earth rotation angle is the defining linear function of the UT1
Julian-day interval $T_u$ from J2000 (Capitaine, Guinot \&
McCarthy~\cite{capitaine2000}; IERS Conventions (2010),
eq.~5.15~\cite{iers2010}):
\begin{equation}
  \mathrm{ERA}(T_u) = 2\pi \bigl( \operatorname{frac}(T_u)
    + 0.7790572732640 + 0.00273781191135448\, T_u \bigr) ,
  \label{eq:frames:era}
\end{equation}
where $\operatorname{frac}(T_u)$ is taken from the two Julian-date parts
separately before the small rate coefficient multiplies the collapsed
$T_u$ --- the arrangement that preserves the two-part precision
(Section~\ref{sec:frames:implementation}). UT1 itself is
$\mathrm{UT1} = \mathrm{TAI} - \Delta AT + \Delta\mathrm{UT1}$ with
$\Delta AT$ from the leap table of Chapter~\ref{ch:time} and constant
$\Delta\mathrm{UT1}$ per assumption 3.

\subsection{Moon principal-axis frame}

The DE-series lunar librations are published as the 3-1-3 Euler angles
$\phi, \theta, \psi$ that carry the ICRF axes onto the lunar
principal axes~\cite{park2021}, so with the sequence convention of
equation~\eqref{eq:rotations:e313},
\begin{equation}
  \dcm{I}{L} = \mat{R}_3(\psi)\, \mat{R}_1(\theta)\, \mat{R}_3(\phi).
  \label{eq:frames:moonpa}
\end{equation}
The construction takes the three angles as plain arguments: the DE440
Chebyshev evaluator that interpolates them is a separate module, and the
assembly must not depend on it. Representative DE440 magnitudes ($\phi$
small, $\theta \approx \qty{0.4}{\radian}$, $\psi$ unbounded) are among
the committed golden angle sets.

\subsection{Mars body-fixed frame}

The IAU Working Group 2015 report~\cite{archinal2018} (with its
correction~\cite{archinal2019}; numerical values as distributed in NAIF
\texttt{pck00011.tpc}~\cite{pck00011}) gives the Mars north-pole right
ascension $\alpha_0$, declination $\delta_0$, and prime-meridian angle
$W$, in degrees, as functions of the TDB interval from J2000.0
($= \mathrm{JD}\,2451545.0$ TDB): $d$ in days for the $W$ rate and $T =
d/36525$ in Julian centuries elsewhere,
\begin{equation}
  \begin{aligned}
    \alpha_0 &= 317.269202 - 0.10927547\,T
      + \textstyle\sum_i a_i \sin(\theta_i), \\
    \delta_0 &= 54.432516 - 0.05827105\,T
      + \textstyle\sum_i b_i \cos(\phi_i), \\
    W &= 176.049863 + 350.891982443297\,d
      + \textstyle\sum_i c_i \sin(\chi_i),
  \end{aligned}
  \label{eq:frames:mars}
\end{equation}
with the periodic terms of Table~\ref{tab:frames:mars}. The body-fixed
frame follows the standard IAU construction:
\begin{equation}
  \dcm{I}{M} = \mat{R}_3(W)\,
    \mat{R}_1\bigl(\tfrac{\pi}{2} - \delta_0\bigr)\,
    \mat{R}_3\bigl(\tfrac{\pi}{2} + \alpha_0\bigr),
  \label{eq:frames:marsdcm}
\end{equation}
whose third row is the pole direction $(\cos\delta_0 \cos\alpha_0,
\cos\delta_0 \sin\alpha_0, \sin\delta_0)$ --- the cross-check the golden
generator asserts independently.

\begin{table}[htbp]
  \centering
  \caption{Mars periodic rotational-element terms (Archinal et
    al.~\cite{archinal2018}; amplitudes and arguments in degrees, argument
    rates in degrees per TDB Julian century).}
  \label{tab:frames:mars}
  \begin{tabular}{lrrr}
    \toprule
    Element (function) & Amplitude & Phase & Rate \\
    \midrule
    $\alpha_0$ ($\sin$) & 0.000068 & 198.991226 & 19139.4819985 \\
                        & 0.000238 & 226.292679 & 38280.8511281 \\
                        & 0.000052 & 249.663391 & 57420.7251593 \\
                        & 0.000009 & 266.183510 & 76560.6367950 \\
                        & 0.419057 & 79.398797 & 0.5042615 \\
    \midrule
    $\delta_0$ ($\cos$) & 0.000051 & 122.433576 & 19139.9407476 \\
                        & 0.000141 & 43.058401 & 38280.8753272 \\
                        & 0.000031 & 57.663379 & 57420.7517205 \\
                        & 0.000005 & 79.476401 & 76560.6495004 \\
                        & 1.591274 & 166.325722 & 0.5042615 \\
    \midrule
    $W$ ($\sin$)        & 0.000145 & 129.071773 & 19140.0328244 \\
                        & 0.000157 & 36.352167 & 38281.0473591 \\
                        & 0.000040 & 56.668646 & 57420.9295360 \\
                        & 0.000001 & 67.364003 & 76560.2552215 \\
                        & 0.000001 & 104.792680 & 95700.4387578 \\
                        & 0.584542 & 95.391654 & 0.5042615 \\
    \bottomrule
  \end{tabular}
\end{table}

\section{Discretization and implementation}
\label{sec:frames:implementation}

\begin{enumerate}
  \item \textbf{Module and traceability.} The chain lives in
    \texttt{cpp/src/frames.cpp}; the source comments echo the labels
    \texttt{eq:frames:fw}, \texttt{eq:frames:nut00b},
    \texttt{eq:frames:xy}, \texttt{eq:frames:s06},
    \texttt{eq:frames:c2i}, \texttt{eq:frames:era},
    \texttt{eq:frames:chain}, \texttt{eq:frames:moonpa}, and
    \texttt{eq:frames:mars} verbatim.
  \item \textbf{Series tables and transcription control.} All series
    coefficients are compiled constants in
    \texttt{cpp/src/frames\_series.hpp}, transcribed verbatim from the
    ERFA 2.0.1 reference source~\cite{erfa201}. Transcription is
    controlled twice: the golden generator carries an independent Python
    mirror of the same tables, validated functionally against ERFA at
    314 epochs at generation time (residuals at the $10^{-16}$ level),
    and the committed table file \texttt{series\_terms.toml} is compared
    term by term against the compiled constants
    (\testid{frames_series_transcription}, exact binary64 equality).
  \item \textbf{Fixed summation order (D-10).} Both Poisson series sum
    smallest terms first --- the reverse of the published table order,
    exactly the SOFA/ERFA evaluation order --- and the polynomial
    combinations use the reference's nested parenthesization, so the
    implementation is bit-compatible with the operation sequence the
    golden values were generated by, up to libm differences.
  \item \textbf{Two-part UT1 through the ERA.} A matrix-element target of
    $10^{-11}$ corresponds to $\sim \qty{1.4e-7}{\second}$ of UT1; a
    single binary64 of seconds since J2000 carries only
    $\sim \qty{2.4e-7}{\second}$ of resolution by 2060
    (Chapter~\ref{ch:time}), which is why the two-part epoch exists. The
    implementation forms UT1 as an exact whole-day part
    ($\mathrm{JD}\,2451544.5 + d$, exact in binary64 over the table
    domain) plus a seconds-of-day fraction $(s - \Delta AT +
    \Delta\mathrm{UT1})/86400$ with quantum
    $\sim \qty{1.5e-11}{\second}$, and
    equation~\eqref{eq:frames:era} takes $\operatorname{frac}(T_u)$ from
    the parts separately; only the $2.7 \times 10^{-3}$ rate coefficient
    multiplies the collapsed $T_u$, suppressing its rounding by the same
    factor. The fraction may be negative for epochs in the first
    $\Delta AT$ seconds of a TAI day; the final modular reduction
    normalizes the angle.
  \item \textbf{One rotation kernel.} Every matrix is composed from the
    $\mat{R}_1/\mat{R}_2/\mat{R}_3$ primitives of
    Chapter~\ref{ch:rotations} (which match the SOFA sign convention), in
    the same association order as the reference routines; the Moon PA
    assembly is literally the kernel's 3-1-3 composition (asserted
    bit-identical in \testid{frames_moon_pa_golden}).
  \item \textbf{Mars argument reduction.} Each periodic argument is
    reduced modulo \qty{360}{\degree} \emph{before} the radian conversion,
    and $W$'s linear term is reduced the same way, in a fixed order
    identical to the golden generator; the TDB day count is formed as the
    exact two-part difference $(\mathrm{jd}_1 - 2451545.0) +
    \mathrm{jd}_2$ before the single collapse.
  \item \textbf{Where polar motion would go.} The omitted $\mat{W}$ of
    the full IERS chain would left-multiply
    equation~\eqref{eq:frames:chain}; the cookbook cross-check test
    documents this by comparing against the published no-polar-motion
    matrix.
\end{enumerate}

\section{Validation evidence}
\label{sec:frames:validation}

Tables~\ref{tab:frames:validation-core}
and~\ref{tab:frames:validation-py} list the tests that constitute this
chapter's validation evidence. Each test identifier is machine-checked to
exist in the test tree by \texttt{scripts/lint\_chapters.py}; golden
provenance and tolerances are recorded in
\texttt{tests/golden/frames/manifest.toml}.

\begin{table}[htbp]
  \centering
  \caption{Validation evidence for the reference-frame module: core
    (doctest, \texttt{cpp/tests/}).}
  \label{tab:frames:validation-core}
  \begin{tabular}{lp{8.6cm}}
    \toprule
    Test ID & Acceptance basis \\
    \midrule
    \testid{frames_series_transcription} &
      Every compiled series constant equals the committed,
      ERFA-validated table term by term (exact binary64 equality). \\
    \testid{frames_erfa_chain_golden} &
      Phase~2 exit criterion 1: $\dcm{I}{E}$ matrix elements match the
      ERFA-generated chain at 14 epochs spanning 2020--2060 to
      $\leq 10^{-11}$ (observed worst \num{3.5e-14}); chain components
      $\Delta\psi, \Delta\epsilon, X, Y, s, \mathrm{ERA}$ to
      $\leq \qty{1e-12}{\radian}$ (observed worst \qty{3.6e-14}{\radian});
      two-part TAI plumbing bit-identical; orthonormality of every
      produced matrix. \\
    \testid{frames_sofa_cookbook_crosscheck} &
      The published SOFA cookbook worked example (2007-04-05 12:00 UTC,
      IAU 2006/2000A, CIO based, no polar motion)~\cite{sofacookbook}:
      matrix elements within the documented \num{8e-9} model-difference
      bound (2000B vs.\ 2000A plus the cookbook's CIP corrections;
      observed \num{2.3e-9}); published $X$, $Y$, $s$, ERA at their
      documented bounds. \\
    \testid{frames_eci_ecef_roundtrip} &
      Phase~2 exit criterion 6b: ECI$\to$ECEF$\to$ECI position round trip
      at $r = \qty{6778137}{\metre}$, 14 epochs $\times$ 6 directions,
      error $\leq \qty{1e-8}{\metre}$ (observed worst
      \qty{3.8e-9}{\metre}); norm preservation to the same bound. \\
    \testid{frames_moon_pa_golden} &
      Equation~\eqref{eq:frames:moonpa} vs.\ ERFA-composed goldens at 7
      committed angle sets (representative DE440 magnitudes plus
      degenerate geometry) to $\leq 10^{-15}$ per element; bit equality
      with the rotation kernel's 3-1-3 composition. \\
    \testid{frames_mars_iau_golden} &
      Equation~\eqref{eq:frames:mars} vs.\ the independent golden
      evaluation at 6 epochs to $\leq \qty{1e-13}{\radian}$ per element
      angle and $\leq 10^{-14}$ per matrix element (observed bit-identical);
      TDB plumbing bit-identical; pole-vector cross-check. \\
    \bottomrule
  \end{tabular}
\end{table}

\begin{table}[htbp]
  \centering
  \caption{Validation evidence for the reference-frame module: binding
    surface (pytest, \texttt{tests/python/}).}
  \label{tab:frames:validation-py}
  \begin{tabular}{lp{7.0cm}}
    \toprule
    Test ID & Acceptance basis \\
    \midrule
    \testid{test_frames_bindings_match_goldens} &
      The pybind11 surface reproduces the ERFA-chain goldens at the same
      $10^{-11}$/$10^{-12}$ gates. \\
    \testid{test_frames_cookbook_bindings} &
      The published cookbook example through the bindings at its
      documented bounds. \\
    \testid{test_frames_eci_ecef_roundtrip_bindings} &
      Criterion 6b through the bindings ($\leq \qty{1e-8}{\metre}$). \\
    \testid{test_frames_moon_mars_bindings} &
      Moon PA and Mars element/matrix goldens through the bindings. \\
    \bottomrule
  \end{tabular}
\end{table}

\section{References}
\label{sec:frames:references}

This chapter relies on the IERS Conventions (2010)~\cite{iers2010} for the
CIO-based transformation structure, the $s$ series, and the ERA
expression; on McCarthy and Luzum~\cite{mccarthy2003} for the IAU 2000B
nutation model and its accuracy statement; on Capitaine, Wallace and
Chapront~\cite{capitaine2003} for the IAU 2006 precession, with the
Fukushima--Williams parameterization per Fukushima~\cite{fukushima2003}
and Wallace and Capitaine~\cite{wallace2006}; on Capitaine, Guinot and
McCarthy~\cite{capitaine2000} for the Earth rotation angle; on the ERFA
reference implementation~\cite{erfa201} for the transcribed series
coefficients and evaluation order; on the SOFA cookbook~\cite{sofacookbook}
for the published worked example; on Park et al.~\cite{park2021} for the
DE440 libration-angle convention; and on Archinal et
al.~\cite{archinal2018, archinal2019} (values per NAIF
\texttt{pck00011.tpc}~\cite{pck00011}) for the Mars rotational elements.
Full bibliographic details appear in the bibliography.
